\documentclass[11pt]{article}

\usepackage[utf8]{inputenc}
\usepackage[T1]{fontenc}
\usepackage{amsmath,amssymb,amsthm}
\usepackage{microtype}
\usepackage[margin=1.05in]{geometry}
\usepackage{booktabs}
\usepackage{array}
\usepackage{xcolor}
\usepackage{listings}
\usepackage[colorlinks=true,linkcolor=blue!50!black,citecolor=blue!50!black,urlcolor=blue!50!black]{hyperref}

\emergencystretch=1.5em

% Unicode characters that appear inside \texttt{} identifiers in prose.
\DeclareUnicodeCharacter{03B4}{\ensuremath{\delta}}   % δ
\DeclareUnicodeCharacter{03B9}{\ensuremath{\iota}}    % ι
\DeclareUnicodeCharacter{03B1}{\ensuremath{\alpha}}   % α
\DeclareUnicodeCharacter{03B2}{\ensuremath{\beta}}    % β
\DeclareUnicodeCharacter{2080}{\ensuremath{_0}}       % ₀
\DeclareUnicodeCharacter{2081}{\ensuremath{_1}}       % ₁
\DeclareUnicodeCharacter{2082}{\ensuremath{_2}}       % ₂
\DeclareUnicodeCharacter{2083}{\ensuremath{_3}}       % ₃

\definecolor{leankw}{RGB}{0,32,96}
\definecolor{leancomment}{RGB}{96,96,96}

\lstdefinelanguage{lean}{
  morekeywords={theorem,lemma,def,noncomputable,abbrev,structure,where,fun,
    match,with,by,intro,exact,refine,obtain,attribute,local,instance,import,
    namespace,end,open,variable,haveI,have,set,let,show,rintro,apply},
  sensitive=true,
  morecomment=[l]{--},
  morecomment=[s]{/-}{-/},
}
\lstset{
  language=lean,
  basicstyle=\small\ttfamily,
  keywordstyle=\color{leankw}\bfseries,
  commentstyle=\color{leancomment}\itshape,
  columns=fullflexible,
  keepspaces=true,
  breaklines=true,
  showstringspaces=false,
  xleftmargin=1em,
  literate=
    {ℤ}{{\ensuremath{\mathbb{Z}}}}1
    {ℕ}{{\ensuremath{\mathbb{N}}}}1
    {ℝ}{{\ensuremath{\mathbb{R}}}}1
    {×}{{\ensuremath{\times}}}1
    {⟶}{{\ensuremath{\longrightarrow}}}2
    {→}{{\ensuremath{\rightarrow}}}1
    {↦}{{\ensuremath{\mapsto}}}1
    {≫}{{\ensuremath{\gg}}}1
    {≅}{{\ensuremath{\cong}}}1
    {⊞}{{\ensuremath{\boxplus}}}1
    {∩}{{\ensuremath{\cap}}}1
    {∪}{{\ensuremath{\cup}}}1
    {∀}{{\ensuremath{\forall}}}1
    {∃}{{\ensuremath{\exists}}}1
    {∈}{{\ensuremath{\in}}}1
    {∉}{{\ensuremath{\notin}}}1
    {⊆}{{\ensuremath{\subseteq}}}1
    {∨}{{\ensuremath{\lor}}}1
    {∧}{{\ensuremath{\land}}}1
    {¬}{{\ensuremath{\neg}}}1
    {≠}{{\ensuremath{\neq}}}1
    {≤}{{\ensuremath{\le}}}1
    {α}{{\ensuremath{\alpha}}}1
    {β}{{\ensuremath{\beta}}}1
    {δ}{{\ensuremath{\delta}}}1
    {ι}{{\ensuremath{\iota}}}1
    {σ}{{\ensuremath{\sigma}}}1
    {λ}{{\ensuremath{\lambda}}}1
    {ε}{{\ensuremath{\varepsilon}}}1
    {φ}{{\ensuremath{\varphi}}}1
    {ψ}{{\ensuremath{\psi}}}1
    {π}{{\ensuremath{\pi}}}1
    {Δ}{{\ensuremath{\Delta}}}1
    {∂}{{\ensuremath{\partial}}}1
    {∑}{{\ensuremath{\textstyle\sum}}}1
    {∐}{{\ensuremath{\coprod}}}1
    {₀}{{\ensuremath{_0}}}1
    {₁}{{\ensuremath{_1}}}1
    {₂}{{\ensuremath{_2}}}1
    {₃}{{\ensuremath{_3}}}1
    {↥}{{\ensuremath{\uparrow}}}1
    {⇑}{{\ensuremath{\Uparrow}}}1
    {•}{{\ensuremath{\cdot}}}1
    {−}{{\ensuremath{-}}}1
    {≃}{{\ensuremath{\simeq}}}1
    {∘}{{\ensuremath{\circ}}}1
    {ᶜ}{{\ensuremath{^{c}}}}1
    {🄯}{{}}1
    {⁻¹}{{\ensuremath{^{-1}}}}2
    {𝟙}{{\ensuremath{\mathbf{1}}}}1
    {ₗ}{{\ensuremath{_l}}}1
    {→₀}{{\ensuremath{\rightarrow_0}}}2
}

\newtheorem{theorem}{Theorem}[section]
\newtheorem{lemma}[theorem]{Lemma}
\newtheorem{proposition}[theorem]{Proposition}
\newtheorem{corollary}[theorem]{Corollary}
\newtheorem{definition}[theorem]{Definition}
\theoremstyle{remark}
\newtheorem{remark}[theorem]{Remark}

\newcommand{\Z}{\mathbb{Z}}
\newcommand{\R}{\mathbb{R}}
\newcommand{\N}{\mathbb{N}}
\newcommand{\leanname}[1]{\texttt{#1}}

\title{Formalizing Singular Homology in Lean 4:\\
Homotopy Invariance, Subdivision, Mayer--Vietoris,\\
and the Homology of Spheres}

\author{Jonathan Washburn\\
\small Recognition Physics Institute, Austin, Texas\\
\small \texttt{washburn@recognitionphysics.org}}

\date{July 18, 2026}

\begin{document}

\sloppy

\maketitle

\begin{abstract}
Mathlib, the mathematical library of the Lean 4 proof assistant, defines
singular homology of topological spaces through the singular simplicial set
and the alternating face map complex, and it proves essentially no theorems
about the resulting functor. This paper reports a formalization, in Lean 4
over Mathlib, of the classical computational layer of singular homology with
integer coefficients: the prism-operator proof of homotopy invariance, the
long exact sequence of a pair, barycentric subdivision with its diameter
estimate, the small-simplices and small-chains theorems, the Mayer--Vietoris
long exact sequence with honest space-level induced maps, the homology of
spheres (vanishing in the intermediate range and nonvanishing in the top
degree), and the corollary that spheres of different dimensions are never
homotopy equivalent. On top of this layer we formalize applications about
complements of embedded circles and arcs in $S^D$ for arbitrary topological
embeddings, with no smoothness, PL, or tameness hypotheses: the complement of
the unknot in $S^3$ has nonvanishing first homology, and the complement of an
embedded arc in $S^D$ has vanishing first homology in every dimension, a
formalized case of Hatcher's Theorem~2B.1. These combine into a client
theorem: an embedded circle in $S^D$ with homologically nontrivial complement
exists exactly when $D = 3$. The development comprises about 8{,}400 lines
across twelve modules, builds with zero \texttt{sorry}, and every capstone
reports the axiom footprint \texttt{[propext, Classical.choice, Quot.sound]}.
\end{abstract}

\section{Introduction}\label{sec:intro}

Singular homology is the first genuinely computable invariant that a student
of algebraic topology meets, and its computational core is remarkably
stable across textbooks: homotopy invariance by the prism operator, the long
exact sequence of a pair, barycentric subdivision and the small-simplices
argument, excision or Mayer--Vietoris, and then the homology of spheres with
its classical corollaries (invariance of dimension, Brouwer's fixed point
theorem, separation theorems). Mathlib~\cite{mathlib}, the library of the
Lean~4 proof assistant~\cite{lean4}, recently acquired the definition:
\leanname{AlgebraicTopology.singularHomologyFunctor} sends a topological
space to the homology of the alternating face map complex of its singular
simplicial set, with coefficients in a module. What Mathlib does not have is
the computational layer. At the time of writing there is essentially no
theorem in Mathlib about the values of this functor: no homotopy invariance,
no exact sequences, no subdivision, no computation of $H_*(S^n)$.

This paper reports a development that builds that layer, working directly
against Mathlib's own definition of singular homology. No project-local
surrogate stands in between. Every
statement below about $H_n(X)$ is a statement about
\leanname{singularHomologyFunctor} applied to the coefficient module $\Z$,
so anything proved here composes with future Mathlib development without
translation. The development goes far enough to support theorems of genuinely
geometric content. Its final applications concern arbitrary topological
embeddings of circles and arcs into spheres, with no smoothness, PL, or
local flatness assumptions: the wild cases are included, and this is a
strength of the homological method that the formalization preserves. As a
single sentence of motivation beyond topology: the client theorem in
Section~\ref{sec:client} (nontrivial linking of embedded circles forces
ambient dimension three) is consumed by a physics program, described in a
companion paper~\cite{companion}, that derives spatial dimensionality from
information-theoretic axioms; nothing in the present paper depends on that
program.

\subsection{Contributions}

All results are formalized in Lean 4 over Mathlib, build with zero
\texttt{sorry}, and are stated for singular homology with $\Z$ coefficients.

\begin{itemize}
\item \textbf{Homotopy invariance} (954 lines). The prism operator, its five
  face identities, the chain homotopy identity
  $\partial P + P\partial = g_\# - f_\#$, and the packaged consequences:
  homotopic maps induce equal maps on homology, and homotopy equivalences
  induce isomorphisms (Hatcher Theorem~2.10, Corollary~2.11).
\item \textbf{The long exact sequence of a pair} (239 lines). For an
  injective continuous map $A \to X$, the short exact sequence of singular
  chain complexes and its homology long exact sequence, with the connecting
  homomorphism (Hatcher Theorem~2.16, for pairs presented by injections).
\item \textbf{Barycentric subdivision} (1{,}587 lines). A reusable
  combinatorial theory of affine chains with cone recursions for the
  subdivision operator $S$ and the homotopy $T$, the identities
  $\partial S = S \partial$ and $\partial T + T\partial = 1 - S$, transport
  to the singular chain complex, the $\tfrac{n}{n+1}$ diameter contraction,
  and the small-simplices theorem \leanname{exists\_sdOpIter\_small}
  (the geometric heart of Hatcher's proof of excision, pp.~119--124).
\item \textbf{Mayer--Vietoris} (1{,}791 lines). The small-chains subcomplex
  $C^{U,V}_*(X)$, the small-chains theorem (the inclusion induces an
  isomorphism on homology; Hatcher Proposition~2.21), the Mayer--Vietoris
  short exact sequence of chain complexes (which needs no openness), and the
  Mayer--Vietoris long exact sequence with induced maps that are honestly the
  space-level maps (openness enters only here).
\item \textbf{Homology of spheres} (745 + 694 lines). The augmentation
  apparatus, $H_0(X) \cong \Z$ for path-connected $X$ (as invertibility of
  the augmentation), vanishing of positive-degree homology for contractible
  spaces, the suspension isomorphism
  $H_{k+2}(S^{n+1}) \cong H_{k+1}(S^{n})$ obtained from the Mayer--Vietoris
  connecting map, and the two dimension detectors: $H_k(S^n) = 0$ for
  $1 \le k \ne n$ and $H_n(S^n) \ne 0$ for $n \ge 1$. Corollary: spheres of
  different dimensions are never homotopy equivalent.
\item \textbf{Complement topology for wild embeddings} (308 + 249 + 629 +
  870 lines). The complement of the flat unknot in $S^3$ has nonvanishing
  $H_1$; no embedded circle in $S^0$ or $S^1$ has homologically nontrivial
  complement; for $D \ge 2$, $D \ne 3$, the circle-complement question
  reduces by Mayer--Vietoris to arcs; and embedded arcs in $S^D$ have
  $H_1$-acyclic complements in every dimension $D$, by a compact-support
  bisection argument (the arc case of Hatcher Theorem~2B.1, in degree one,
  for arbitrary topological embeddings).
\item \textbf{The client theorem} (311 + 42 lines). The detector
  \leanname{DetectsNontrivialLinking}~$D$ (some embedded circle in $S^D$
  has complement with nonvanishing $H_1$) holds if and only if $D = 3$;
  packaged as a content-typed binder with an anti-cheat gate.
\end{itemize}

Table~\ref{tab:contrib} maps modules to their classical references. All
twelve modules live in the directory
\path{IndisputableMonolith/Foundation/} of the source repository; a public
mirror is available (Section~\ref{sec:artifact}).

\begin{table}[t]
\centering
\small
\begin{tabular}{@{}lll@{}}
\toprule
Module & Content & Classical reference \\
\midrule
\leanname{SingularPrism} & prism operator, homotopy invariance & Hatcher Thm.~2.10, Cor.~2.11 \\
\leanname{SingularPair} & chain SES, LES of a pair & Hatcher Thm.~2.16 \\
\leanname{SingularSubdivision} & barycentric subdivision, small simplices & Hatcher pp.~119--124 \\
\leanname{SingularMayerVietoris} & small chains, MV sequence & Hatcher Prop.~2.21, \S 2.2 \\
\leanname{SingularSphere} & $H_0$, contractible vanishing, MV tools & Hatcher Prop.~2.6, Cor.~2.11 \\
\leanname{SingularSphereGeometry} & sphere homology, dimension detectors & Hatcher Cor.~2.14 \\
\leanname{UnknotComplementRetract} & unknot complement, $H_1 \ne 0$ & retraction argument \\
\leanname{LinkingVanishingLowDim} & no detection in $S^0$, $S^1$ & point-set topology \\
\leanname{LinkingVanishingHighDim} & reduction of circles to arcs & Hatcher \S 2.2, \S 2.B \\
\leanname{ArcComplementAcyclic} & arc complements $H_1$-acyclic, all $D$ & Hatcher Thm.~2B.1 (arc case) \\
\leanname{PublicSpine} & detector, binder, anti-cheat gate & --- \\
\leanname{PublicSpineLinkingClosure} & \leanname{forces\_D3}, \leanname{target\_D3} & --- \\
\bottomrule
\end{tabular}
\caption{The twelve modules and their classical anchors. Line counts by
\texttt{wc -l}: 954, 239, 1587, 1791, 745, 694, 308, 249, 629, 870, 311, 42
(total 8{,}419).}
\label{tab:contrib}
\end{table}

\subsection{Related work}

Singular homology has been formalized before, in other systems and with
other designs. Harrison developed singular homology in HOL
Light~\cite{harrison-homology}, reaching the Eilenberg--Steenrod axioms and
consequences including Brouwer's fixed point theorem and Jordan--Brouwer
separation; Paulson ported that development (about 11{,}400 lines) to
Isabelle/HOL, where it ships with the distribution as
\texttt{HOL-Homology}~\cite{paulson-porting}. Those developments use a
simple-type-theoretic encoding of spaces and hand-rolled free abelian
groups. In the dependently typed world, the homotopy type theory community
computes homotopy and cohomology groups synthetically: Brunerie's
$\pi_4(S^3) \simeq \Z/2$~\cite{brunerie}, cellular cohomology in
HoTT~\cite{buchholtz-favonia}, and machine computation of the Brunerie
number in cubical Agda~\cite{ljungstrom-mortberg}. Synthetic results speak
about types with path structure; the singular homology of general topological
spaces studied here is a different object, so the two do not directly overlap. In
Lean, Murphy formalized the Brouwer fixed point theorem via a singular
homology apparatus built for that purpose, before Mathlib's
\leanname{singularHomologyFunctor} existed~\cite{birs-cohomology}. Mathlib
itself has extensive simplicial machinery (simplicial sets, Dold--Kan,
simplicial homotopy), group and Galois cohomology, and de Rham cohomology in
the geometry library, but the singular homology functor arrived bare. The
present development differs from all of the above in one design decision:
it proves theorems about Mathlib's own functor, in Mathlib's category
language (\leanname{TopCat}, \leanname{ModuleCat} $\Z$, homological complexes),
so that the layer can be upstreamed statement by statement
(Section~\ref{sec:conclusion}).

\section{Background and the target object}\label{sec:background}

\subsection{Singular homology in Mathlib}

Mathlib defines the singular simplicial set of a space $X$ as the restricted
Yoneda presheaf of the standard simplices: an $n$-simplex is a continuous map
$\Delta^n \to X$, where $\Delta^n$ is \leanname{stdSimplex} $\R$
\leanname{(Fin (n+1))}. Applying the free-module functor levelwise and taking
the alternating face map complex produces the singular chain complex functor
\leanname{singularChainComplexFunctor}; its homology is
\leanname{singularHomologyFunctor}. The development fixes coefficients once
and for all:

\begin{lstlisting}
noncomputable abbrev SC (X : TopCat.{0}) : ChainComplex (ModuleCat.{0} ℤ) ℕ :=
  ((AlgebraicTopology.singularChainComplexFunctor (ModuleCat.{0} ℤ)).obj
    (ModuleCat.of ℤ ℤ)).obj X

noncomputable abbrev Hgrp (X : TopCat.{0}) (n : ℕ) : ModuleCat.{0} ℤ :=
  (SC X).homology n
\end{lstlisting}

Both are reducible abbreviations, so $\leanname{Hgrp}\ X\ n$ \emph{is} the
value of Mathlib's singular homology functor at $X$ with $\Z$ coefficients;
there is no project-local homology theory to reconcile. Throughout the paper
$H_n(X)$ abbreviates this object and $C_n(X)$ the degree-$n$ chain group.

A recurring working presentation: with $\Z$ coefficients the degree-$n$
chain group is the coproduct $\coprod_{\sigma} \Z$ in
$\leanname{ModuleCat}\ \Z$, indexed by the singular $n$-simplices. The
development names the index type \leanname{Idx}~$X$~$n$, the coproduct
\leanname{Cgrp}, and the generator inclusion \leanname{gen}~$X$~$n$~$\sigma$.
Two normal-form lemmas (\leanname{gen\_d} for the boundary,
\leanname{gen\_map} for functoriality) reduce every chain-level computation
to signed sums over generators. This presentation is set up once in
\leanname{SingularPrism} and reused by every later module.

\subsection{Design choices}

The development is classical. Choice and excluded middle are used freely, as
is customary in Mathlib's topology and category theory; there is no attempt
at constructive content, and none of the results are claimed constructively.
There is also no synthetic homotopy theory: spaces are topological spaces,
homotopies are continuous maps $I \times X \to Y$, and the results are about
the classical singular functor. Universe level is pinned to
\leanname{TopCat.\{0\}}, which is where Mathlib's spheres live.

Every capstone theorem in the development reports the axiom footprint
\begin{center}
\texttt{[propext, Classical.choice, Quot.sound]}
\end{center}
under \texttt{\#print axioms}: the three standard Lean axioms, and nothing
else. No result depends on any project-local axiom.

\section{Homotopy invariance: the prism operator}\label{sec:prism}

The module \leanname{SingularPrism} (954 lines) formalizes the textbook
prism-operator proof of homotopy invariance, following Hatcher's proof of
Theorem~2.10 but reorganized so that the combinatorial cancellation is a
standalone abstract lemma.

\subsection{Topological prisms and face identities}

The prism decomposition writes $\Delta^n \times I$ as a union of $(n+1)$
simplices of dimension $n+1$. Rather than triangulate literally, the
formalization works with the $(n+1)$ affine maps
$\leanname{prism}\ i : \Delta^{n+1} \to \Delta^n \times I$ for
$i : \leanname{Fin}\ (n+1)$, and proves the five face identities that the
cancellation argument needs: the composites of $\leanname{prism}\ i$ with
the topological face inclusions
$\leanname{face}\ j : \Delta^n \to \Delta^{n+1}$ satisfy
\leanname{prism\_comp\_face\_top}, \leanname{prism\_comp\_face\_bot},
\leanname{prism\_comp\_face\_cancel}, \leanname{prism\_comp\_face\_of\_le},
and \leanname{prism\_comp\_face\_of\_gt}. These are proved by direct
computation with the affine coordinates of \leanname{stdSimplex}.

Given a homotopy $H : C(I \times X, Y)$ and a singular $n$-simplex $\sigma$
of $X$, each prism map produces a singular $(n+1)$-simplex of $Y$
(\leanname{prismSimplex}), and the prism operator is the signed sum over the
decomposition, extended to chains by the universal property of the
coproduct:

\begin{lstlisting}
noncomputable def prismOp {X Y : TopCat.{0}} (H : C(I × X, Y)) (n : ℕ)
    (prisms : Fin (n + 1) →
      C(stdSimplex ℝ (Fin (n + 2)), stdSimplex ℝ (Fin (n + 1)) × I)) :
    Cgrp X n ⟶ Cgrp Y (n + 1) :=
  Sigma.desc (Pgen H n prisms)
\end{lstlisting}

\subsection{The chain homotopy}

The heart of the proof is an abstract double-sum cancellation
(\leanname{prism\_sum\_cancellation}): for coefficients in any abelian
group, the alternating double sum over faces of prisms collapses to the top
and bottom faces once the diagonal terms cancel in pairs. The five face
identities feed exactly this lemma, transported through the generator normal
forms. The result is the chain homotopy identity, stated in positive degrees
and in degree zero:

\begin{theorem}[prism chain homotopy;
\leanname{prism\_chain\_homotopy\_succ}, \leanname{\dots\_zero}]
\label{thm:prism}
Let $F_0, F_1 : X \to Y$ be continuous and $H$ a homotopy from $F_0$ to
$F_1$. Writing $P$ for the prism operator of $H$ and $f_\#$ for the induced
chain map, on degree $n+1$ chains
\[
\partial \circ P + P \circ \partial = (F_1)_\# - (F_0)_\#,
\qquad\text{and on degree } 0,\quad
\partial \circ P = (F_1)_\# - (F_0)_\#.
\]
\end{theorem}

\begin{lstlisting}
lemma prism_chain_homotopy_succ (Ho : ContinuousMap.Homotopy F₀.hom F₁.hom) (n : ℕ) :
    bnd X n ≫ prismOp Ho.toContinuousMap n prism +
        prismOp Ho.toContinuousMap (n + 1) prism ≫ bnd Y (n + 1) =
      chainMap F₁ (n + 1) - chainMap F₀ (n + 1)
\end{lstlisting}

Packaging the two identities as a Mathlib \leanname{Homotopy} of chain
complex morphisms yields the module's exports.

\begin{theorem}[homotopy invariance;
\leanname{homotopic\_maps\_induce\_same\_homology}]\label{thm:hinv}
Homotopic continuous maps $f, g : X \to Y$ induce equal maps
$H_n(X) \to H_n(Y)$ for every $n$, where $H_n$ is Mathlib's
\leanname{singularHomologyFunctor} with $\Z$ coefficients.
\end{theorem}

\begin{theorem}[homotopy equivalence;
\leanname{homotopyEquiv\_homology\_iso}]\label{thm:heq}
A homotopy equivalence $X \simeq Y$ induces an isomorphism
$H_n(X) \cong H_n(Y)$ for every $n$, and the map induced by the forward
direction is an isomorphism in $\leanname{ModuleCat}\ \Z$
(\leanname{isIso\_homology\_map\_of\_homotopyEquiv}).
\end{theorem}

These are Hatcher Theorem~2.10 and Corollary~2.11. Every downstream module
consumes homotopy invariance only through Theorem~\ref{thm:heq}.

\section{Relative homology and the long exact sequence of a pair}
\label{sec:pair}

The module \leanname{SingularPair} (239 lines) is short because it delegates
the homological algebra to Mathlib's \leanname{ShortComplex.ShortExact} API.
The geometric input is one lemma: for an \emph{injective} continuous map
$f : A \to X$, the induced chain map is a split monomorphism in every degree
(\leanname{chainMap\_mono}, \leanname{sChainMap\_mono}). Injectivity
suffices, no embedding hypothesis is needed, because a singular simplex of
$X$ that factors through the image factors uniquely through $A$ at the level
of index types, and the retraction \leanname{genRetract} is defined on
generators by that factorization.

The relative complex is a cokernel, and the pair sequence is short exact:

\begin{lstlisting}
noncomputable def relSC (f : A ⟶ X) : ChainComplex (ModuleCat.{0} ℤ) ℕ :=
  cokernel (sChainMap f)

lemma pairSES_shortExact (f : A ⟶ X) (hf : Function.Injective f.hom) :
    (pairSES f).ShortExact
\end{lstlisting}

\begin{theorem}[LES of a pair; \leanname{pairδ},
\leanname{pair\_les\_exact₁}, \leanname{pair\_les\_exact₂},
\leanname{pair\_les\_exact₃}]\label{thm:les}
For injective continuous $f : A \to X$ there is a connecting homomorphism
$\partial : H_{n+1}(X, A) \to H_n(A)$, and the sequence
\[
\cdots \to H_{n+1}(X) \to H_{n+1}(X, A) \xrightarrow{\ \partial\ }
H_n(A) \to H_n(X) \to H_n(X, A) \to \cdots
\]
is exact at every position, in all degrees including zero.
\end{theorem}

This is Hatcher Theorem~2.16 for pairs presented by injections. A sanity
theorem closes the module: $H_n(X, X) = 0$
(\leanname{relative\_homology\_id\_isZero}). The same split-mono lemma is
reused in Section~\ref{sec:mv} for the Mayer--Vietoris columns.

\section{Barycentric subdivision and small simplices}\label{sec:subdivision}

The module \leanname{SingularSubdivision} (1{,}587 lines) is the technical
center of the development. The excision-grade input that Mayer--Vietoris
needs is: any singular chain can be subdivided, without changing its
homology class, until every simplex lands in $U$ or in $V$. Textbooks prove
this with about four pages of formulas (Hatcher pp.~119--124); the
formalization's main design decision is to do all the algebra on a
combinatorial layer where the formulas become structural recursions.

\subsection{The affine chain layer}

An affine $n$-simplex in a type $\alpha$ is just a vertex tuple, and an
affine chain is a finitely supported $\Z$-linear combination of tuples:

\begin{lstlisting}
abbrev AC (α : Type) (n : ℕ) : Type := (Fin (n + 1) → α) →₀ ℤ
\end{lstlisting}

Faces are tuple reindexings along \leanname{Fin.succAbove}; the cone with
apex $b$ is \leanname{Fin.cons}~$b$; the boundary \leanname{abnd} is the
alternating sum of faces; the augmentation \leanname{eps} sums coefficients.
No topology appears at this layer, and, importantly, no geometry either: the
barycenter enters as an \emph{abstract apex function}
$\leanname{bary} : (\leanname{Fin}\ (m+1) \to \alpha) \to \alpha$, a
parameter of the constructions. The subdivision operator and its chain
homotopy are the classical cone recursions
\[
S\lambda = b_\lambda \cdot S(\partial\lambda),
\qquad
T\lambda = b_\lambda \cdot (\lambda - T(\partial\lambda)),
\]
implemented as \leanname{asub} and \leanname{atee} by structural recursion
on degree. All chain identities hold at this level for \emph{any} apex
function (degree zero needs only $\leanname{bary}\ 0\ w = w\ 0$): the cone
boundary identity, $\partial\partial = 0$, the chain-map property
$\partial S = S\partial$ (\leanname{abnd\_comp\_asub}), and the homotopy
identity:

\begin{theorem}[affine chain homotopy; \leanname{abnd\_comp\_atee}]
\label{thm:atee}
For every apex function \leanname{bary} and every degree $n$,
$\partial \circ T + T \circ \partial = \mathrm{id} - S$ on affine
$(n+1)$-chains, where $S = \leanname{asub}$ and $T = \leanname{atee}$.
\end{theorem}

\begin{lstlisting}
theorem abnd_comp_atee (bary : ∀ {m : ℕ}, (Fin (m + 1) → α) → α) :
    ∀ n, (abnd (n + 1)).comp (atee bary (n + 1)) + (atee bary n).comp (abnd n)
      = LinearMap.id - asub bary (n + 1)
\end{lstlisting}

The iterates $S^k$ come with a telescoped homotopy
(\leanname{ateeIter}, \leanname{abnd\_comp\_ateeIter}):
$\partial T_k + T_k \partial = 1 - S^k$. Proving these identities by
structural recursion over a two-line grammar of generators is where the
formalization gains the most over a literal transcription of the textbook,
which manipulates the same content through nested summation indices.

\subsection{Geometry: the honest barycenter and the diameter estimate}

Stage 5b instantiates the abstract layer on
$\alpha = \leanname{stdSimplex}\ \R\ (\leanname{Fin}\ (d+1))$: the
realization \leanname{affineMap} of a vertex tuple as a continuous affine
map, and the honest barycenter \leanname{sbary}, with its equivariance
\leanname{sbary\_affineMap} under affine maps. Equivariance is what lets
the abstract identities transport: pushforward \leanname{amap} along a
barycenter-preserving map commutes with \leanname{asub} and
\leanname{atee}.

The quantitative content is Stage 6. The lemma
\leanname{asub\_support\_bound} tracks supports through the cone recursion:
each piece of a subdivided simplex has vertex spread at most
$\tfrac{n}{n+1}$ times the original, by the convex-hull maximum principle.
Iterating
(\leanname{asubIter\_support\_bound}) gives the geometric decay
$\bigl(\tfrac{n}{n+1}\bigr)^k$, and therefore:

\begin{theorem}[uniform smallness; \leanname{exists\_asubIter\_small}]
\label{thm:small-affine}
For every $\varepsilon > 0$ and every vertex tuple $w$ on the standard
simplex, there is an iterate $k$ such that every tuple in the support of
$S^k(w)$ has all pairwise vertex distances below $\varepsilon$.
\end{theorem}

\subsection{Transport and the small-simplices theorem}

Stage 5c transports the operators to the singular chain complex. On a
generator $\sigma$, the singular subdivision \leanname{sdOp} is the
pushforward along $\sigma$ of the subdivision of the identity tuple; the
homotopy \leanname{tOp} is transported the same way. The chain-map property
(\leanname{sdOp\_comp\_bnd}), the homotopy identities
(\leanname{tOp\_chain\_homotopy\_succ}, \leanname{tOp\_chain\_homotopy\_zero}),
naturality in the space (\leanname{sdOp\_natural}, \leanname{tOp\_natural}),
and the iterated versions (\leanname{sdOpIter}, \leanname{tOpIter},
\leanname{tOpIter\_chain\_homotopy\_succ}) all descend from the affine layer
through the generator normal forms.

The capstone combines Theorem~\ref{thm:small-affine} with the Lebesgue
number lemma on the compact metric space $\Delta^n$:

\begin{theorem}[small-simplices theorem; \leanname{exists\_sdOpIter\_small}]
\label{thm:small-simplices}
Let $U, V$ be open with $U \cup V = X$ and let $s$ be a singular
$n$-simplex of $X$. Then there is an iterate $k$ such that every piece of
$\leanname{sdOpIter}\ X\ n\ k$ applied to the generator of $s$ has range
contained in $U$ or contained in $V$.
\end{theorem}

\begin{lstlisting}
theorem exists_sdOpIter_small {n : ℕ} (U V : Set X) (hU : IsOpen U) (hV : IsOpen V)
    (hUV : U ∪ V = Set.univ) (s : Idx X n) :
    ∃ k, ∀ u ∈ ((asubIter (baryFn n) k n) (asimplex (idTuple n))).support,
      Set.range ⇑(simplexEquiv X n (pushSimplex (simplexEquiv X n s) u)) ⊆ U ∨
        Set.range ⇑(simplexEquiv X n (pushSimplex (simplexEquiv X n s) u)) ⊆ V
\end{lstlisting}

Together with the homotopy witness $1 - S^k = \partial T_k + T_k \partial$,
this is exactly the input the Mayer--Vietoris layer consumes: subdivision
reaches the small subcomplex, and the reached chain is homologous to the
original.

\section{Mayer--Vietoris}\label{sec:mv}

The module \leanname{SingularMayerVietoris} (1{,}791 lines) assembles the
subdivision layer into the Mayer--Vietoris long exact sequence. Fix
$U, V \subseteq X$ (arbitrary subsets at first).

\subsection{The small-chains subcomplex and the small-chains theorem}

A singular simplex is \leanname{Small}~$U$~$V$ if its range lies in $U$ or
in $V$; the small chains $C^{U,V}_n(X)$ are the span of the small
generators, assembled into a subcomplex \leanname{SSC}~$U$~$V$ with
inclusion \leanname{smallι}. The main theorem of the first half is
Hatcher's Proposition~2.21 specialized to two-set covers:

\begin{theorem}[small-chains theorem; \leanname{smallι\_isIso\_homologyMap},
\leanname{smallChainsHomologyIso}]\label{thm:small-chains}
For open $U, V$ with $U \cup V = X$, the inclusion
$C^{U,V}_*(X) \to C_*(X)$ induces an isomorphism on homology in every
degree.
\end{theorem}

The proof is elementwise, through a concrete surjectivity-and-injectivity
criterion for homology maps of complexes of $\Z$-modules
(\leanname{isIso\_homologyMap\_of\_elementwise}): every cycle of $C_*(X)$
subdivides into the subcomplex (Theorem~\ref{thm:small-simplices}, with the
iterate $k$ uniformized over the finitely many generators of the chain by
taking the maximum), and the homotopy operator $T_k$, which preserves the
subcomplex by naturality, witnesses that nothing is lost. Hatcher proves
more (a chain homotopy inverse); the homology-level isomorphism is what
Mayer--Vietoris needs, and it is what the formalization exports.

\subsection{The Mayer--Vietoris sequence}

The short exact sequence lives at the chain level and needs no openness at
all:

\begin{lstlisting}
noncomputable def mvα :
    SC (TopCat.of (U ∩ V : Set X)) ⟶ SC (TopCat.of U) ⊞ SC (TopCat.of V) :=
  biprod.lift (sChainMap (mvInclU U V)) (-(sChainMap (mvInclV U V)))

noncomputable def mvβ :
    SC (TopCat.of U) ⊞ SC (TopCat.of V) ⟶ SSC U V :=
  biprod.desc (smallU U V) (smallV U V)

theorem mvSES_shortExact : (mvSES U V).ShortExact
\end{lstlisting}

The observation that
$0 \to C_*(U \cap V) \to C_*(U) \oplus C_*(V) \to C^{U,V}_*(X) \to 0$
is exact for \emph{arbitrary} subsets, openness entering only when the small
subcomplex is compared to the full complex, is standard but easy to blur in
informal treatments; the formal statement keeps the hypotheses honest.
Degreewise exactness in the middle
(\leanname{mv\_middle\_exact}) is a direct combinatorial argument on
coproduct coordinates: a pair of chains agreeing after inclusion must come
from the intersection, generator by generator.

Combining with Theorem~\ref{thm:small-chains} and the homology LES of a
short exact sequence of complexes gives the space-level sequence. The maps
matter: \leanname{mvPair} is (induced by) the two inclusions
$U \cap V \to U$, $U \cap V \to V$ with a sign, \leanname{mvSum} is the sum
of the maps induced by $U \to X$, $V \to X$, and the connecting map
\leanname{mvδ} is the LES connecting map conjugated by the small-chains
isomorphism. All three are honest space-level maps of Mathlib's homology
functor, so downstream users can compose them with arbitrary induced maps.

\begin{theorem}[Mayer--Vietoris LES; \leanname{mvδ},
\leanname{mv\_exact₁}, \leanname{mv\_exact₂}, \leanname{mv\_exact₃},
\leanname{mvSum\_epi\_zero}]\label{thm:mv}
For open $U, V$ with $U \cup V = X$, the sequence
\[
\cdots \to H_{n+1}(X) \xrightarrow{\ \partial\ } H_n(U \cap V)
\to H_n(U) \oplus H_n(V) \to H_n(X) \to \cdots
\]
is exact at every position; in degree zero the final map is an epimorphism.
\end{theorem}

\begin{remark}[engineering: the $\leanname{ModuleCat}\ \Z$ instance diamond]
\label{rem:diamond}
One engineering lesson deserves a highlighted place, because it will affect
anyone who does elementwise arguments in complexes of
$\leanname{ModuleCat}\ \Z$. The elements of a chain group carry a
$\Z$-module structure through two unification paths: the categorical
instance of $\leanname{ModuleCat}\ \Z$, and the generic chain
$\leanname{AddCommGroup.toIntModule}$ (with its
$\leanname{SubNegMonoid.toZSMul}$ scalar action). The two are propositionally
but not definitionally equal, and Lean's elaborator will happily pick
different paths in different subterms of one goal, producing goals that look
identical and refuse to close. Sixteen hundred lines of
\leanname{Finsupp} and span manipulation became feasible only after pinning
the priorities locally:
\begin{lstlisting}
attribute [local instance 0] AddCommGroup.toIntModule
attribute [local instance 0] SubNegMonoid.toZSMul
\end{lstlisting}
Demoting the two generic instances to the lowest priority makes elaboration
choose the categorical instance uniformly. The same two lines are repeated
in every downstream module that computes with elements
(\leanname{SingularSphere}, \leanname{SingularSphereGeometry},
\leanname{LinkingVanishingHighDim}, \leanname{ArcComplementAcyclic}).
\end{remark}

\section{The homology of spheres}\label{sec:spheres}

Two modules compute sphere homology: \leanname{SingularSphere} (745 lines)
builds the abstract tools, \leanname{SingularSphereGeometry} (694 lines)
supplies the geometry. The sphere model is
$\leanname{Sph}\ n = \leanname{Metric.sphere}\ (0 : \R^{n+1})\ 1$ as an
object of \leanname{TopCat}.

\subsection{Degree zero and contractible spaces}

The augmentation \leanname{augTo} sends every $0$-simplex to
$1 \in \Z$; on homology it gives $\leanname{augH} : H_0(X) \to \Z$, and a
choice of point gives a section $\leanname{ptH}$.

\begin{theorem}[$H_0$ of a path-connected space;
\leanname{isIso\_augH\_of\_pathConnected}]\label{thm:h0}
For a path-connected space $X$, the augmentation
$H_0(X) \to \Z$ is an isomorphism in $\leanname{ModuleCat}\ \Z$.
\end{theorem}

This is Hatcher Proposition~2.6, stated as invertibility of a specific map so
that later arguments can chase the augmentation through induced maps; a bare
isomorphism statement would not support that chase. Positive degrees of a point vanish by
inspection of the complex
(\leanname{isZero\_homology\_of\_totallyDisconnected} handles any totally
disconnected space), and homotopy invariance transports this:

\begin{theorem}[contractible vanishing;
\leanname{isZero\_homology\_of\_contractible}]\label{thm:contractible}
For contractible $X$ and $m \ne 0$, the object $H_m(X)$ is a zero object.
\end{theorem}

The module also proves the abstract Mayer--Vietoris induction tools: if
$U, V$ are open, cover $X$, and have contractible-like vanishing, the
connecting map is an isomorphism in the relevant range
(\leanname{isIso\_mvδ}, \leanname{isIso\_mvδ\_of\_contractible}), plus the
degree-one edge cases (\leanname{isZero\_h1\_of\_contractible},
\leanname{mono\_mvPair\_zero}).

\subsection{Geometry: covers, stereographic projection, and suspension}

The geometric module covers $S^{n+1}$ by the complements
$U = S^{n+1} \setminus \{\text{south}\}$ and
$V = S^{n+1} \setminus \{\text{north}\}$. Each is contractible: the
formalization proves
\leanname{contractibleSpace\_compl\_singleton\_sphere} through Mathlib's
stereographic projection. The intersection deformation-retracts to the
equator: \leanname{interHomeoPunctured} identifies $U \cap V$ with a
punctured hyperplane, polar coordinates \leanname{puncturedPolar} split it
as $S^n \times (0,\infty)$-like data, and
\leanname{hequivProdContractible} collapses the contractible factor,
producing the homotopy equivalence \leanname{interHomotopyEquiv} between
$U \cap V$ and $S^n$. Feeding these into the abstract Mayer--Vietoris
induction yields the suspension isomorphism as a composite of the
connecting map and the equator identification:

\begin{lstlisting}
noncomputable def suspensionIso (n k : ℕ) :
    Hgrp (Sph (n + 1)) (k + 2) ≅ Hgrp (Sph n) (k + 1)
\end{lstlisting}

\subsection{The dimension detectors}

Induction on $n$ with base cases $S^0$ (two points; homology vanishes in
positive degrees by total disconnectedness) and $S^1$ (where the degree-one
case is handled by exactness with both cover pieces contractible) gives the
two exports that all applications consume:

\begin{theorem}[vanishing; \leanname{sphere\_homology\_vanish}]
\label{thm:vanish}
For all $n$ and all $k$ with $1 \le k$ and $k \ne n$, the object
$H_k(S^n)$ is a zero object.
\end{theorem}

\begin{theorem}[nonvanishing; \leanname{h1\_s1\_ne\_zero},
\leanname{sphere\_top\_ne\_zero}]\label{thm:nonvanish}
$H_1(S^1)$ is not a zero object, and for every $n \ge 1$, $H_n(S^n)$ is
not a zero object.
\end{theorem}

A remark on honesty of scope: the development exports vanishing and
nonvanishing, and does not export an isomorphism $H_n(S^n) \cong \Z$ for
general $n$. Nonvanishing is proved by chasing a nonzero connecting map down
the suspension ladder to $H_1(S^1) \ne 0$; the exactness bookkeeping for
that chase does not require identifying the group. For $n = 1$ specifically,
a separate, earlier module of the same repository
(\leanname{CircleWindingChain}, 7{,}753 lines, predating this campaign)
proves the isomorphism
$H_1(S^1) \cong \Z$ as \leanname{circleH1ZIsoInt\_holds}, by an explicit
winding-chain computation; Section~\ref{sec:applications} uses it. Deriving
$H_n(S^n) \cong \Z$ for all $n$ from \leanname{suspensionIso} and the
$n = 1$ case is a short induction left for the upstreaming pass
(Section~\ref{sec:conclusion}).

The classical corollary follows by comparing detectors across a homotopy
equivalence:

\begin{theorem}[spheres detect dimension;
\leanname{spheres\_not\_homotopyEquivalent}]\label{thm:spheres-detect}
For $m \ne n$ there is no homotopy equivalence between $S^m$ and $S^n$;
equivalently, a homotopy equivalence between $S^m$ and $S^n$ forces
$m = n$ (\leanname{sphere\_dim\_eq\_of\_homotopyEquiv}).
\end{theorem}

\begin{lstlisting}
theorem spheres_not_homotopyEquivalent {m n : ℕ} (hmn : m ≠ n) :
    IsEmpty (ContinuousMap.HomotopyEquiv ↥(Sph m) ↥(Sph n))
\end{lstlisting}

\section{Applications: complements of wild circles and arcs}
\label{sec:applications}

The application layer studies, for each ambient dimension $D$, whether some
embedded circle in $S^D$ has a complement with nonvanishing first homology.
Everything in this section works with arbitrary
\leanname{Topology.IsEmbedding} maps. No smoothness, PL structure, local
flatness, or tameness is assumed anywhere, so wild embeddings (Antoine,
Fox--Artin, and worse) are within scope of every statement. This generality
costs nothing extra, because the method is homological throughout; it is a
concrete payoff of formalizing the singular theory itself, which a tame
surrogate would not deliver. The scope is deliberately narrow in a different direction: the
development proves a corridor of Alexander-type complement control
(complements of arcs in degree one, plus the specific circle statements
below), and does not prove, or claim, full Alexander duality.

\subsection{Dimension three: the unknot has nontrivial complement}

The module \leanname{UnknotComplementRetract} (308 lines) proves existence
at $D = 3$ with a self-contained retraction argument. The flat unknot is
$z \mapsto (z, 0, 0)$; the core circle of its complement is
$w \mapsto (0, 0, w)$; and the retraction normalizes the last two
coordinates:

\begin{theorem}[unknot complement; \leanname{unknotComplementH1\_ne\_zero}]
\label{thm:unknot}
Assume $H_1(S^1) \cong \Z$. Then the first singular homology of the
complement of the flat unknot in $S^3$ is not a zero object.
\end{theorem}

\begin{lstlisting}
theorem unknotComplementH1_ne_zero
    (h1 : Nonempty ((((AlgebraicTopology.singularHomologyFunctor (ModuleCat ℤ) 1).obj
      (ModuleCat.of ℤ ℤ)).obj (TopCat.sphere.{0} 1)) ≅ ModuleCat.of ℤ ℤ)) :
    ¬ CategoryTheory.Limits.IsZero
      (((AlgebraicTopology.singularHomologyFunctor (ModuleCat ℤ) 1).obj
        (ModuleCat.of ℤ ℤ)).obj
        (TopCat.of {x : TopCat.sphere.{0} 3 // x ∉ Set.range unknot}))
\end{lstlisting}

Functoriality alone finishes: \leanname{retract\_core} shows the composite
$S^1 \to \mathrm{complement} \to S^1$ is the identity, so $H_1(S^1) \cong \Z$
is a retract of $H_1(\mathrm{complement})$, which therefore cannot be zero.
The hypothesis is discharged by \leanname{circleH1ZIsoInt\_holds}
(Section~\ref{sec:spheres}) when the statement is consumed in
Section~\ref{sec:client}. Notably, this existence half uses no homotopy
invariance, no subdivision, and no Mayer--Vietoris: it is a 308-line
functoriality argument available even before the computational layer.

\subsection{Low dimensions: no room for linking}

The module \leanname{LinkingVanishingLowDim} (249 lines) disposes of
$D = 0$ and $D = 1$ by point-set topology.

\begin{theorem}[low-dimensional vanishing; \leanname{not\_detects\_zero},
\leanname{not\_detects\_one}]\label{thm:lowdim}
No embedded circle in $S^0$ or in $S^1$ has a complement with nonvanishing
$H_1$.
\end{theorem}

For $D = 0$ there is no embedded circle at all: $S^0$ is a two-point
discrete space and homology of any subspace is concentrated in degree zero,
by \leanname{isZero\_homology\_of\_totallyDisconnected} again. For $D = 1$,
a continuous injection from $S^1$ to itself is surjective, by
\leanname{continuous\_injective\_circle\_self\_surjective}. That proof
runs through the intermediate value theorem, via
\leanname{no\_continuous\_injective\_circle\_to\_real}. So the complement
of an embedded circle in $S^1$ is empty and its homology vanishes.

\subsection{High dimensions: reduction to arcs}

The module \leanname{LinkingVanishingHighDim} (629 lines) handles all
$D \ge 2$, $D \ne 3$ at once, conditionally on one clean frontier
hypothesis:

\begin{lstlisting}
def ArcComplementsAcyclic (D : ℕ) : Prop :=
  ∀ a : C(unitInterval, ↥(Sph D)), Topology.IsEmbedding a →
    CategoryTheory.Limits.IsZero
      (Hgrp (TopCat.of {y : ↥(Sph D) // y ∉ Set.range a}) 1)
\end{lstlisting}

The reduction is Mayer--Vietoris. Split the embedded circle into two closed
semicircular arcs meeting in two points. The complements of the two arcs
form an open cover of the two-point complement
$S^D \setminus \{p, q\}$, their intersection is the circle complement, and
the sequence
\[
H_2\bigl(S^D \setminus \{p,q\}\bigr) \to H_1(\text{circle complement})
\to H_1(\text{arc complement}) \oplus H_1(\text{arc complement})
\]
traps the middle term. The two outer terms vanish: the arcs by hypothesis,
and the left term because $S^D \setminus \{p, q\}$ is homotopy equivalent to
$S^{D-1}$ (\leanname{twoPointComplHEquiv}, built from stereographic
projection, a translation homeomorphism, and the polar splitting), whose
$H_2$ vanishes for $D \ne 3$ by Theorem~\ref{thm:vanish}
(\leanname{isZero\_h2\_twoPointCompl}). A pleasant surprise fell out of the
uniform treatment: the case $D = 2$ needs no Jordan--Schoenflies theorem and
no separation theory whatsoever, because the argument only ever consults
$H_2(S^1)$, which vanishes. The classical intuition that dimension two is
special (complements of circles in $S^2$ disconnect) is real but
irrelevant to this statement.

\begin{theorem}[high-dimensional reduction;
\leanname{forces\_D3\_of\_arcAcyclic}]\label{thm:highdim}
Fix $D$ with $1 \le D$ and $D \ne 3$, and assume
$\leanname{ArcComplementsAcyclic}\ D$. Then every embedded circle in $S^D$
has $H_1$-acyclic complement
(\leanname{isZero\_h1\_complement\_of\_embedding}). Consequently, granting the arc hypothesis for
all $D \ge 2$, $D \ne 3$: if some embedded circle in $S^D$ has
homologically nontrivial complement, then $D = 3$.
\end{theorem}

\subsection{Arc complements are acyclic in every dimension}

The module \leanname{ArcComplementAcyclic} (870 lines) discharges the
hypothesis unconditionally, in all dimensions including $D = 3$. This is the
arc case of Hatcher's Theorem~2B.1, in degree one, for arbitrary topological
embeddings.

\begin{theorem}[arc complements; \leanname{arcComplementsAcyclic}]
\label{thm:arc}
For every $D$ and every topological embedding
$a : [0,1] \to S^D$, the object
$H_1\bigl(S^D \setminus a([0,1])\bigr)$ is a zero object.
\end{theorem}

The proof is the classical bisection. Suppose some $1$-cycle $z$ in the
complement of the arc is not a boundary. The module first builds an
elementwise toolkit for homology classes of complexes of $\Z$-modules
(\leanname{classOf}, \leanname{classOf\_eq\_zero\_iff},
\leanname{classOf\_natural}) and its topological wrappers
(\leanname{cls}, \leanname{exists\_nonbounding},
\leanname{bounds\_of\_isZero}). Splitting the arc at its midpoint, the
complements of the two half-arcs cover the complement of their common point;
Mayer--Vietoris exactness at $H_1(U \cap V)$
(\leanname{bounds\_of\_mv}), with $H_2$ and $H_1$ of the punctured sphere
vanishing by Section~\ref{sec:spheres} (the punctured sphere is
contractible), shows $z$ must fail to bound in one of the two half-arc
complements (\leanname{bounds\_of\_halves}). Iterating produces a nested
sequence of parameter intervals (\leanname{badSeq}) converging to a single
point $t^*$. The complement of the point $a(t^*)$ is contractible, so $z$
bounds there; a bounding chain is a finite object, its compact support
avoids $a(J)$ for a small enough interval $J \ni t^*$ around the limit
(\leanname{exists\_chain\_lift}, the compact-support lifting lemma), and
that contradicts badness of the interval $J$. The limit step is where
formalization pressure was highest: ``a chain lives in a finite subcomplex''
is a triviality on paper and a genuine lemma
(\leanname{cPush}, \leanname{cLift}, \leanname{exists\_chain\_lift}) about
filtered behavior of \leanname{Finsupp}-supported chains here.

Theorem~\ref{thm:arc} states degree one only, which is exactly what
Theorem~\ref{thm:highdim} consumes. The full Theorem~2B.1 (all reduced
degrees, and disk complements in all codimensions) remains open in this
development; the bisection scaffolding is degree-generic and should extend.

\section{The client theorem}\label{sec:client}

The module \leanname{PublicSpine} (311 lines) states the campaign-facing
definitions. The 42-line module \leanname{PublicSpineLinkingClosure}
closes them. The detector is content-typed: it is defined on Mathlib's homology
object itself, so no proxy re-statement can satisfy it.

\begin{lstlisting}
def DetectsNontrivialLinking (D : ℕ) : Prop :=
  ∃ f : C(TopCat.sphere.{0} 1, TopCat.sphere.{0} D),
    Topology.IsEmbedding f ∧
      ¬ CategoryTheory.Limits.IsZero (linkingComplementH1 D f)
\end{lstlisting}

\begin{theorem}[the client theorem, in two halves;
\leanname{detectsNontrivialLinking\_three}, \leanname{forces\_D3}]
\label{thm:client}
Existence at three: $\leanname{DetectsNontrivialLinking}\ 3$ holds
(\leanname{detectsNontrivialLinking\_three}), which is Theorem~\ref{thm:unknot}
with the hypothesis discharged by \leanname{circleH1ZIsoInt\_holds}.
Uniqueness: for every $D$, $\leanname{DetectsNontrivialLinking}\ D$ implies
$D = 3$ (\leanname{forces\_D3}), assembled from Theorems~\ref{thm:lowdim},
\ref{thm:highdim} and \ref{thm:arc}. Together the two halves give
$\leanname{DetectsNontrivialLinking}\ D$ if and only if $D = 3$. No single Lean
theorem states that biconditional; it is the conjunction of the two named
declarations.
\end{theorem}

The two halves are packaged for the consuming program by a separate bridge
declaration \leanname{target\_D3}, described below.

\begin{lstlisting}
theorem forces_D3 :
    ∀ D, PublicSpine.DetectsNontrivialLinking D → D = 3 :=
  PublicSpineLinkingAssembly.forces_D3_of_arcAcyclic
    (fun D _ _ => ArcComplementAcyclic.arcComplementsAcyclic D)
\end{lstlisting}

One paragraph on the packaging, with details in the companion
paper~\cite{companion}. The consuming program requires the result
as a binder, \leanname{AlexanderLinkingBridge}, a structure whose fields are
the $H_1(S^1) \cong \Z$ isomorphism, the detection statement at $D = 3$, and
the uniqueness implication; \leanname{target\_D3} inhabits it
unconditionally. Because autonomous provers contributed to this campaign,
the repository enforces an anti-cheat discipline. A CI gate script inspects
the source file \texttt{PublicSpine.lean}: it checks that the target stays a
\leanname{def}, rejects a \texttt{theorem target\_D3} declaration in that
file, and flags known drift patterns in the guidepost \texttt{Skeleton}
module. The gate is a source-text check on those files and does not scan the
whole repository. The binder fields are content-typed against Mathlib's own
objects (a homology object of a concrete subspace of a concrete sphere), so
the target cannot be ``satisfied'' by redefining the statement, weakening the
detector, or introducing a parallel homology theory. The conditional assembly
(\leanname{PublicSpineLinkingAssembly}, 49 lines) was kept as a separate
module so that the exact logical distance between the Mayer--Vietoris
reduction and the unconditional closure stayed visible while the arc lemma
was still open.

\section{Artifact and reproducibility}\label{sec:artifact}

\subsection{Repository and build}

The development lives in the \texttt{IndisputableMonolith} Lean library;
the twelve modules of this paper are under
\texttt{IndisputableMonolith/Foundation/}. A public mirror of the campaign
modules is available at
\begin{center}
\url{https://github.com/jonwashburn/shape-of-logic},
commit \texttt{3606692}\,(\texttt{80a2a84f28b67f2a0edd99756fba7c9fc}).
\end{center}
The library builds with Lake against a pinned Mathlib. Each module builds as
a single target, for example:

\begin{lstlisting}[language={}]
$ lake build IndisputableMonolith.Foundation.SingularMayerVietoris
$ lake build IndisputableMonolith.Foundation.PublicSpineLinkingClosure
\end{lstlisting}

All twelve modules build green with zero \texttt{sorry}. The axiom audit on
the capstones (\leanname{arcComplementsAcyclic}, \leanname{forces\_D3},
\leanname{target\_D3}, and the exports of every section above) reports
exactly:

\begin{lstlisting}[language={}]
#print axioms IndisputableMonolith.Foundation.PublicSpineLinkingClosure.target_D3
-- 'IndisputableMonolith.Foundation.PublicSpineLinkingClosure.target_D3'
--   depends on axioms: [propext, Classical.choice, Quot.sound]
\end{lstlisting}

\subsection{Theorem index}

Table~\ref{tab:index} maps every numbered statement of this paper to its
Lean declaration and module.

\begin{table}[ht]
\centering
\small
\begin{tabular}{@{}lll@{}}
\toprule
Paper statement & Lean declaration(s) & Module \\
\midrule
Thm.~\ref{thm:prism} & \leanname{prism\_chain\_homotopy\_succ}, \leanname{\dots\_zero} & \leanname{SingularPrism} \\
Thm.~\ref{thm:hinv} & \leanname{homotopic\_maps\_induce\_same\_homology} & \leanname{SingularPrism} \\
Thm.~\ref{thm:heq} & \leanname{homotopyEquiv\_homology\_iso}, & \leanname{SingularPrism} \\
 & \leanname{isIso\_homology\_map\_of\_homotopyEquiv} & \\
Thm.~\ref{thm:les} & \leanname{pairδ}, \leanname{pair\_les\_exact₁}/\leanname{₂}/\leanname{₃} & \leanname{SingularPair} \\
$H_n(X,X)=0$ & \leanname{relative\_homology\_id\_isZero} & \leanname{SingularPair} \\
Thm.~\ref{thm:atee} & \leanname{abnd\_comp\_atee}, \leanname{abnd\_comp\_ateeIter} & \leanname{SingularSubdivision} \\
Thm.~\ref{thm:small-affine} & \leanname{asubIter\_support\_bound}, & \leanname{SingularSubdivision} \\
 & \leanname{exists\_asubIter\_small} & \\
Thm.~\ref{thm:small-simplices} & \leanname{exists\_sdOpIter\_small} & \leanname{SingularSubdivision} \\
Thm.~\ref{thm:small-chains} & \leanname{smallι\_isIso\_homologyMap}, & \leanname{SingularMayerVietoris} \\
 & \leanname{smallChainsHomologyIso} & \\
MV chain SES & \leanname{mvSES\_shortExact} & \leanname{SingularMayerVietoris} \\
Thm.~\ref{thm:mv} & \leanname{mvδ}, \leanname{mv\_exact₁}/\leanname{₂}/\leanname{₃}, \leanname{mvSum\_epi\_zero} & \leanname{SingularMayerVietoris} \\
Thm.~\ref{thm:h0} & \leanname{isIso\_augH\_of\_pathConnected} & \leanname{SingularSphere} \\
Thm.~\ref{thm:contractible} & \leanname{isZero\_homology\_of\_contractible} & \leanname{SingularSphere} \\
Suspension iso & \leanname{suspensionIso} & \leanname{SingularSphereGeometry} \\
Thm.~\ref{thm:vanish} & \leanname{sphere\_homology\_vanish} & \leanname{SingularSphereGeometry} \\
Thm.~\ref{thm:nonvanish} & \leanname{h1\_s1\_ne\_zero}, \leanname{sphere\_top\_ne\_zero} & \leanname{SingularSphereGeometry} \\
Thm.~\ref{thm:spheres-detect} & \leanname{spheres\_not\_homotopyEquivalent}, & \leanname{SingularSphereGeometry} \\
 & \leanname{sphere\_dim\_eq\_of\_homotopyEquiv} & \\
Thm.~\ref{thm:unknot} & \leanname{unknotComplementH1\_ne\_zero} & \leanname{UnknotComplementRetract} \\
Thm.~\ref{thm:lowdim} & \leanname{not\_detects\_zero}, \leanname{not\_detects\_one} & \leanname{LinkingVanishingLowDim} \\
Thm.~\ref{thm:highdim} & \leanname{isZero\_h1\_complement\_of\_embedding}, & \leanname{LinkingVanishingHighDim} \\
 & \leanname{not\_detects\_of\_arcAcyclic}, & \\
 & \leanname{forces\_D3\_of\_arcAcyclic} & \\
Thm.~\ref{thm:arc} & \leanname{arcComplementsAcyclic} & \leanname{ArcComplementAcyclic} \\
Detector & \leanname{DetectsNontrivialLinking} & \leanname{PublicSpine} \\
Thm.~\ref{thm:client} ($\Leftarrow$) & \leanname{detectsNontrivialLinking\_three} & \leanname{PublicSpine} \\
Thm.~\ref{thm:client} ($\Rightarrow$) & \leanname{forces\_D3}, \leanname{target\_D3} & \leanname{PublicSpineLinkingClosure} \\
$H_1(S^1)\cong\Z$ (imported) & \leanname{circleH1ZIsoInt\_holds} & \leanname{CircleWindingChain} \\
\bottomrule
\end{tabular}
\caption{Theorem index. All declarations are in namespace
\leanname{IndisputableMonolith.Foundation}, in the sub-namespace matching
the module name.}
\label{tab:index}
\end{table}

\section{Conclusion and upstreaming}\label{sec:conclusion}

The classical computational layer of singular homology now exists in Lean 4
over Mathlib's own definitions: homotopy invariance, the LES of a pair,
barycentric subdivision with its metric estimate, small chains,
Mayer--Vietoris, sphere homology detectors, and a set of complement theorems
strong enough to characterize an ambient dimension. Three formalization
lessons stand out. First, doing the subdivision algebra on an abstract
affine-chain layer with a parametric apex turned the hardest textbook pages
into structural recursions, and the geometry entered exactly once, through
the barycenter's equivariance and the diameter bound. Second, elementwise
arguments in $\leanname{ModuleCat}\ \Z$ complexes are entirely feasible at
this scale, but only after the instance diamond of
Remark~\ref{rem:diamond} is pinned. Third, hypotheses stayed more honest
than the textbook default: injectivity (not embedding) drives the pair
sequence, openness is absent from the Mayer--Vietoris short exact sequence
and enters only at the small-chains comparison, and compact support became
a named lemma where the textbook offers only a shrug.

Several layers are natural Mathlib upstreaming candidates, in dependency
order: the generator normal forms for the singular chain complex
(\leanname{Idx}/\leanname{gen}/\leanname{gen\_d}/\leanname{gen\_map}); the
prism operator and Theorems~\ref{thm:hinv} and \ref{thm:heq}, which are
stated against Mathlib declarations already; the affine-chain subdivision
theory, which is self-contained combinatorics; the small-chains theorem and
the Mayer--Vietoris sequence; and the augmentation apparatus with the sphere
detectors. No part of this development has been merged into Mathlib, and no
claim of that kind is made here; adaptation to Mathlib's evolving simplicial
API, and generalization from $\Z$ to arbitrary coefficients where the proofs
allow it, is the main anticipated work. Beyond upstreaming, two mathematical
continuations are within reach of the existing scaffolding: the isomorphism
$H_n(S^n) \cong \Z$ by induction from the winding-chain base case through
\leanname{suspensionIso}, and the extension of Theorem~\ref{thm:arc} to all
degrees and to disk complements, which would move the development from a
corridor of Alexander-type control toward duality proper.

\subsection*{Acknowledgments}

The author thanks the Lean and Mathlib communities for the foundations this
work builds on.

\begin{thebibliography}{99}

\bibitem{companion}
J.~Washburn.
\newblock \emph{The Unique Dimension in Which an Embedded Circle Links: A
  Machine-Checked Proof that Complement $H_1$ Forces $D = 3$}.
\newblock In preparation, 2026.

\bibitem{hatcher}
A.~Hatcher.
\newblock \emph{Algebraic Topology}.
\newblock Cambridge University Press, 2002.

\bibitem{mathlib}
The mathlib Community.
\newblock The Lean mathematical library.
\newblock In \emph{Proceedings of the 9th ACM SIGPLAN International
  Conference on Certified Programs and Proofs (CPP 2020)}, pages 367--381.
  ACM, 2020.

\bibitem{lean4}
L.~de~Moura and S.~Ullrich.
\newblock The Lean 4 theorem prover and programming language.
\newblock In \emph{Automated Deduction -- CADE 28}, LNCS 12699, pages
  625--635. Springer, 2021.

\bibitem{harrison-homology}
J.~Harrison.
\newblock A HOL Light formalization of singular homology.
\newblock Presentation at Hales60: From the Fundamental Lemma to Discrete
  Geometry, to Formal Verification, 2018.

\bibitem{paulson-porting}
L.~C. Paulson.
\newblock Porting libraries of mathematics between proof assistants.
\newblock Blog post and the \texttt{HOL-Homology} session of the Isabelle
  distribution, 2022.

\bibitem{brunerie}
G.~Brunerie.
\newblock \emph{On the homotopy groups of spheres in homotopy type theory}.
\newblock PhD thesis, Universit\'e Nice Sophia Antipolis, 2016.

\bibitem{buchholtz-favonia}
U.~Buchholtz and K.~Hou~(Favonia).
\newblock Cellular cohomology in homotopy type theory.
\newblock In \emph{Proceedings of the 33rd Annual ACM/IEEE Symposium on
  Logic in Computer Science (LICS 2018)}, pages 521--529. ACM, 2018.

\bibitem{ljungstrom-mortberg}
A.~Ljungstr\"om and A.~M\"ortberg.
\newblock Formalizing $\pi_4(S^3) \cong \mathbb{Z}/2\mathbb{Z}$ and
  computing a Brunerie number in cubical Agda.
\newblock In \emph{Proceedings of the 38th Annual ACM/IEEE Symposium on
  Logic in Computer Science (LICS 2023)}. IEEE, 2023.

\bibitem{birs-cohomology}
A.~Topaz et al.
\newblock Formalization of cohomology theories.
\newblock BIRS workshop report 23w5124, Banff International Research
  Station, 2023.

\end{thebibliography}

\end{document}
